% Method overview.
% (a) QA pruner: one token set for all layers, frames emptied -> timestamps lost.
% (b) Coverage: every layer keeps its own set (top-k of that layer's question attention) under an equal per-frame
%     quota, so every timestamp survives in every layer.
% (c) PCD at one time value of the teaser example (illustrative probabilities): greedy decoding follows the
%     no-video prior to "0"; the contrast promotes "64", which only the memory supports.
\begin{figure*}[t]
\centering
\resizebox{\linewidth}{!}{%
\begin{tikzpicture}[
  font=\sffamily\small,
  tok/.style={minimum size=3.2mm, inner sep=0pt, draw=black!25, fill=black!5},
  keep/.style={tok, fill=ours!80, draw=ours},
  ts/.style={rounded corners=1pt, inner sep=1.2pt, font=\sffamily\scriptsize, fill=orange!25, draw=orange!70},
  tsoff/.style={ts, fill=black!4, draw=black!20, text=black!30},
  panel/.style={rounded corners=4pt, draw=black!35, inner sep=6pt},
  box/.style={rounded corners=3pt, draw=black!45, align=center, inner sep=4pt},
  arr/.style={-{Stealth[length=2.4mm]}, thick, black!60},
  title/.style={font=\sffamily\footnotesize\bfseries, anchor=west},
  note/.style={font=\sffamily\scriptsize, anchor=north west, align=left}
]
%% =============================================================== (a) QA pruner
\node[title] at (0, 6.25) {(a) QA pruner: one token set, shared by all layers};
\foreach \f/\t in {0/0s,1/8s,2/16s,3/24s,4/32s,5/40s} {
  \pgfmathsetmacro{\x}{\f*1.25}
  \ifnum\f=1 \node[tsoff] at (\x+0.32, 5.8) {\t};
  \else\ifnum\f=4 \node[tsoff] at (\x+0.32, 5.8) {\t};
  \else \node[ts] at (\x+0.32, 5.8) {\t}; \fi\fi
  \foreach \i in {0,1,2} \foreach \j in {0,1,2} \node[tok] at (\x+\i*0.32, 5.4-\j*0.32) {};
}
\foreach \f/\i/\j in {0/0/0,0/1/1,2/0/0,2/1/0,2/2/1,2/1/2,3/0/1,3/2/2,5/2/0,5/1/1,5/0/2,5/2/2} {
  \pgfmathsetmacro{\x}{\f*1.25}  \node[keep] at (\x+\i*0.32, 5.4-\j*0.32) {};
}
\node[note, text=loss] at (0, 4.62) {8\,s and 32\,s keep no token: their timestamps are lost};

%% =============================================================== (b) Coverage, three layers
\node[title] at (0, 3.95) {(b) Coverage: equal quota per frame, own set per layer};
% layer slices drawn back to front with an offset
\foreach \L/\dy/\dx/\pat in {36/0.0/0.42/2, 18/-1.1/0.21/1, 1/-2.2/0.0/0} {
  \begin{scope}[shift={(\dx, \dy)}]
    \fill[white, rounded corners=2pt, draw=black!20] (-0.12, 3.62) rectangle (7.3, 2.62);
    \node[font=\sffamily\scriptsize, anchor=east, text=black!60] at (-0.15, 3.12) {$\ell{=}\L$};
    \foreach \f/\t in {0/0s,1/8s,2/16s,3/24s,4/32s,5/40s} {
      \pgfmathsetmacro{\x}{\f*1.25}
      \node[ts, font=\sffamily\tiny] at (\x+0.32, 3.47) {\t};
      \foreach \i in {0,1,2} \node[tok] at (\x+\i*0.32, 3.12) {};
      \foreach \i in {0,1,2} \node[tok] at (\x+\i*0.32, 2.80) {};
    }
    % pattern 0/1/2: different kept tokens per layer, always 2 per frame
    \ifnum\pat=0
      \foreach \f/\i/\j/\k/\l in {0/0/0/2/1,1/1/1/2/0,2/2/0/0/1,3/1/0/1/1,4/0/1/2/1,5/2/0/1/1}
        {\pgfmathsetmacro{\x}{\f*1.25} \node[keep] at (\x+\i*0.32, 3.12-\j*0.32) {}; \node[keep] at (\x+\k*0.32, 3.12-\l*0.32) {};}
    \fi
    \ifnum\pat=1
      \foreach \f/\i/\j/\k/\l in {0/1/1/2/0,1/0/0/2/1,2/1/0/1/1,3/0/1/2/1,4/1/0/2/0,5/0/0/0/1}
        {\pgfmathsetmacro{\x}{\f*1.25} \node[keep] at (\x+\i*0.32, 3.12-\j*0.32) {}; \node[keep] at (\x+\k*0.32, 3.12-\l*0.32) {};}
    \fi
    \ifnum\pat=2
      \foreach \f/\i/\j/\k/\l in {0/2/0/1/1,1/0/1/1/0,2/0/0/2/1,3/2/0/0/1,4/0/0/1/1,5/1/0/2/1}
        {\pgfmathsetmacro{\x}{\f*1.25} \node[keep] at (\x+\i*0.32, 3.12-\j*0.32) {}; \node[keep] at (\x+\k*0.32, 3.12-\l*0.32) {};}
    \fi
  \end{scope}
}
\node[note, text=black!70] at (0, 0.28)
  {layer $\ell$ keeps $\operatorname{top}_{B/F}$ of the question's attention $a_\ell(i)$ in each frame:\\
   every timestamp survives; each layer picks its own evidence};

\node[panel, fit={(-0.6,6.5) (8.0,-0.55)}] (mem) {};

%% =============================================================== (c) PCD
\node[title] at (9.0, 6.25) {(c) PCD at one time value: ``Seen: \underline{\hspace{1.2em}} seconds''};
\begin{axis}[
    name=pm, at={(9.0cm, 3.55cm)}, anchor=south west, width=4.3cm, height=3.1cm,
    ybar, bar width=9pt, ymin=0, ymax=0.95, symbolic x coords={0,64,103}, xtick=data,
    xticklabel style={font=\sffamily\scriptsize}, ytick=\empty, axis x line*=bottom, axis y line=none,
    title={$p(y\mid\mathcal{M},x)$: pruned memory}, title style={font=\sffamily\scriptsize, yshift=-4pt},
    nodes near coords, nodes near coords style={font=\sffamily\tiny, /pgf/number format/fixed, /pgf/number format/precision=2}, enlarge x limits=0.3]
  \addplot[fill=ours!70, draw=ours] coordinates {(0,0.38) (64,0.34) (103,0.12)};
\end{axis}
\begin{axis}[
    name=p0, at={(9.0cm, 0.15cm)}, anchor=south west, width=4.3cm, height=3.1cm,
    ybar, bar width=9pt, ymin=0, ymax=0.95, symbolic x coords={0,64,103}, xtick=data,
    xticklabel style={font=\sffamily\scriptsize}, ytick=\empty, axis x line*=bottom, axis y line=none,
    title={$p(y\mid\varnothing,x)$: no video (prior)}, title style={font=\sffamily\scriptsize, yshift=-4pt},
    nodes near coords, nodes near coords style={font=\sffamily\tiny, /pgf/number format/fixed, /pgf/number format/precision=2}, enlarge x limits=0.3]
  \addplot[fill=black!30, draw=black!50] coordinates {(0,0.80) (64,0.05) (103,0.05)};
\end{axis}
\begin{axis}[
    name=sc, at={(14.6cm, 1.55cm)}, anchor=south west, width=4.3cm, height=3.1cm,
    ybar, bar width=9pt, ymin=0, ymax=0.95, symbolic x coords={0,64,103}, xtick=data,
    xticklabel style={font=\sffamily\scriptsize}, ytick=\empty, axis x line*=bottom, axis y line=none,
    title={$\operatorname{softmax}_y s_t(y)$: contrasted}, title style={font=\sffamily\scriptsize, yshift=-4pt},
    nodes near coords, nodes near coords style={font=\sffamily\tiny, /pgf/number format/fixed,
    /pgf/number format/precision=2}, enlarge x limits=0.3]
  \addplot[fill=orange!55, draw=orange!80!black] coordinates {(0,0.15) (64,0.72) (103,0.13)};
\end{axis}
\node[box, fill=orange!8, font=\sffamily\scriptsize, text width=4.5cm, anchor=north west] at (14.4, 6.05)
  {$s_t(y)=(1{+}\alpha_t)\log p_{\mathcal{M}}(y)-\alpha_t\log p_{\varnothing}(y)$\\[2pt]
   $\alpha_t = 1-\max_y p_{\mathcal{M}}(y) = 0.62$; only on time values};
\draw[arr] ($(pm.east)+(0.1,0)$) -- ++(0.75,-0.6);
\draw[arr] ($(p0.east)+(0.1,0.3)$) -- ++(0.75,0.6);
\node[font=\sffamily\scriptsize, text=loss, anchor=west, align=left] at (9.05, -0.45) {greedy: argmax $p_{\mathcal{M}}$ $\rightarrow$ ``0'' (the prior)};
\node[font=\sffamily\scriptsize\bfseries, text=gain, anchor=west, align=left] at (14.6, 0.85) {\method: ``64'' (only the memory supports it)};
\node[box, fill=white, font=\ttfamily\scriptsize, align=left, anchor=north west, text width=8.9cm] (out) at (9.0, -0.8)
  {Seen: 65.0 seconds, a little boy is seen on the train\\
   Seen: 94.0 seconds, a lady is seen talking to the little boy\\
   Answer: ... for about 29 seconds.\quad Time: [65.0, 94.0]\hfill{\sffamily\color{black!55}$\leftarrow$ answer reserve}};
\draw[arr] (mem.east |- 0,4.2) -- node[below, font=\sffamily\scriptsize] {pruned KV} (9.0, 4.2);
\end{tikzpicture}}
\caption{\textbf{Overview of \method.} \textbf{(a)}~QA-oriented pruners keep one set of salient or unique tokens in
every layer and can empty whole frames, removing their timestamps. \textbf{(b)}~Coverage prunes each layer
separately: every frame keeps the same share of the layer budget, filled with the tokens that the question attends to
most \emph{in that layer}. Every timestamp survives in every layer, and different layers keep different evidence
(sharing one set across layers costs $1.5$ mIoU, \cref{tab:shared}). \textbf{(c)}~PCD at one time value of the
teaser example (probabilities illustrative). Greedy decoding picks ``0'', which the memory supports only as much as
the video-free prior does; the contrasted score discounts it and promotes ``64'', which only the memory supports.
The answer reserve guarantees that the final interval is written (output shown is the real \method answer).}
\label{fig:overview}
\end{figure*}


\section{Method: \method}
\label{sec:method}

\method intervenes at the two places where \cref{sec:diagnosis} found the failures. Coverage acts on the memory and
prior-contrastive decoding (PCD) acts on the decoder, and both sit behind a frame-selection front end (BOLT). No
component needs training, and \cref{fig:overview} gives an overview.

\subsection{Coverage: equal quotas filled by the question}
\label{sec:coverage}

Once the $F$ frames are encoded, we prune the memory a single time, separately in each of the $L$ layers. Let
$\mathcal{T}_f$ be the visual tokens of frame pair $f \in \{1,\dots,F'\}$ and $B$ the per-layer budget for visual
tokens.

\paragraph{Question relevance.}
We append the question $x$ to the full, uncompressed memory and run it exactly through all layers, so the queries
at layer $\ell$ come from the real hidden states. A visual token $i$ is scored at layer $\ell$ by the attention it
receives from the question, averaged over heads $h$ and question positions $j$:
\begin{equation}
  a_\ell(i) \;=\; \frac{1}{H\,|x|} \sum_{h=1}^{H} \sum_{j \in x}
  \operatorname{softmax}_i\!\Big( \tfrac{1}{\sqrt{d}}\, Q^{(\ell,h)}_j \, {K^{(\ell,h)}_i}^{\!\top} \Big).
  \label{eq:relevance}
\end{equation}

\paragraph{Equal quota.}
Each frame pair gets the same share of the budget, $B_f = \lfloor B/F' \rfloor$, with any remainder handed out in
order. Layer $\ell$ then keeps
\begin{equation}
  \mathcal{K}_\ell \;=\; \bigcup_{f=1}^{F'} \topop_{B_f} \big\{\, a_\ell(i) \;:\; i \in \mathcal{T}_f \,\big\}.
  \label{eq:quota}
\end{equation}
The timestamp tokens of $f$ are kept if and only if $\mathcal{K}_\ell \cap \mathcal{T}_f \neq \emptyset$. At any
practical budget the equal quota leaves every frame pair with at least one token, so every time anchor survives,
while the question decides which part of each frame is shown. Kept tokens keep their original rotary positions,
which preserves their temporal order.

\paragraph{Why equal, and why the question.}
We chose the equal quota as a deliberately simple coverage constraint. It guarantees the timeline, which grounding
needs and multiple-choice QA does not. It also keeps attention from collapsing onto a few frames. When quotas follow
attention mass instead, attention sinks and recency bias can give a single frame up to 227 of the roughly 2{,}140
kept tokens. In a pre-registered test, relevance-weighted quotas with a floor, computed from SigLIP similarity or
from attention mass, did not beat equal quotas ($-0.2$ mIoU).
Attention becomes informative once it comes from the \emph{actual} question. Streaming memories such as
HERMES~\cite{hermes} score tokens with generic probe questions instead, and generic probes with a global top-$k$ are
the weakest variant we measured (18.3 mIoU); adding the equal quota lifts them by $5.3$. With the actual question,
however, a global top-$k$ without any quota has so far done as well as the equal quota (\cref{sec:res-coverage}).
In that setting the quota acts as a guarantee and is not where the gain comes from.

\paragraph{Per-layer selection.}
Each layer keeps its own set $\mathcal{K}_\ell$. If we force a single shared set by averaging attention over
layers, \method changes by $-1.5$ mIoU, while random selection is unaffected ($+0.7$). Per-layer selection
therefore helps question-driven scores specifically, and it does not explain why random selection beats input-level
pruners.

\paragraph{Frames and budget.}
An equal quota only works if each frame gets enough tokens to show something. With 2--4 tokens per frame pair, the
most-attended tokens carry little content and Coverage collapses (\cref{sec:frames}). For small budgets we therefore
spend the tokens on fewer frames. At a fixed number of kept tokens, 32 frames with twice the share per frame recover
what 64 frames lose. Our rule is to use $F{=}64$ when each frame pair keeps at least ${\sim}8$ tokens and to halve
$F$ otherwise. Kept timestamps count against the budget, and at budgets of a few hundred tokens they take up a large
part of it, so some of the collapse may come from this accounting (\cref{sec:budget}).
\TODO{this rule was derived from the results of \cref{sec:frames}; validate it on held-out data and with the fixed
timestamp rule.}

\paragraph{Cost.}
For each question, Coverage adds one forward pass of the question over the uncompressed memory. Because the full
memory is prefilled first, Coverage shrinks the memory that the decoder reads but does not lower the cost of
building it. PCD decodes a second, prompt-only row in lockstep at every step, even though its scores are used only
inside time values. This roughly doubles the decoder's per-step work, on a much shorter context. \TODO{measured prefill, time to first token, total latency, peak
memory and KV size; the current implementation runs 3--4 scoring passes, reducible to one.}

\subsection{PCD: prior-contrastive decoding of time values}
\label{sec:contrast}

Let $p_\Mem(y) = p(y \mid \Mem, x, y_{<t})$ be the next-token distribution with the compressed memory and
$p_\varnothing(y) = p(y \mid \varnothing, x, y_{<t})$ the distribution with the same prompt and generated prefix but
no visual tokens at all. Then
\begin{equation}
  \log p_\Mem(y) \;=\; \log p_\varnothing(y) \;+\; \underbrace{\log \frac{p_\Mem(y)}{p_\varnothing(y)}}_{\text{evidence from } \Mem},
  \label{eq:decomp}
\end{equation}
so when the memory says nothing about $y$, the evidence term is flat and the video-free prior decides. This is the
failure described in \cref{sec:prior}. Prior-contrastive decoding (PCD) amplifies the evidence term,
\begin{equation}
\begin{aligned}
  s_t(y) &\;=\; (1+\alpha_t)\log p_\Mem(y) - \alpha_t \log p_\varnothing(y) \\
         &\;=\; \log p_\varnothing(y) + (1+\alpha_t) \log \frac{p_\Mem(y)}{p_\varnothing(y)},
\end{aligned}
  \label{eq:pcd}
\end{equation}
and picks $y_t = \argmax_{y \in \mathcal{V}_\beta} s_t(y)$ over the plausible set
$\mathcal{V}_\beta = \{ y : p_\Mem(y) \geq \beta \max_{y'} p_\Mem(y') \}$ with $\beta{=}0.1$, as in VCD~\cite{vcd},
so that tokens the memory makes implausible are never promoted. We adapt this rule to grounding in three ways.

\paragraph{No-video counterfactual.}
We use an empty negative memory instead of a noised or shuffled one. This isolates exactly the prior measured in
\cref{sec:prior}, and it is cheap, since its cache holds only the prompt.

\paragraph{Time scope.}
PCD is active only while the model writes a time value, that is, inside a \texttt{Seen:} time before its
description and inside \texttt{Time: [$s$, $e$]}. Every other token, including the answer text, is decoded greedily
from $p_\Mem$, so PCD leaves the answer alone wherever the memory already provides it.

\paragraph{Confidence-adaptive strength.}
We set $\alpha_t = \alpha_{\max}\,\big(1 - \max_y p_\Mem(y)\big)$ with $\alpha_{\max}{=}1$. When the memory makes the
model confident, \eg because the evidence really starts at 0\,s, PCD fades out and the early answer stands. When the
model hesitates between several times, the evidence term decides.

The prompt format matters here. Under the timeline prompt, PCD first acts on the listed \texttt{Seen:} moments and
then on the interval, so the interval is written from observations that have already been contrasted. With PCD we
also reserve 64 tokens of the generation budget for the final \texttt{Answer:} and \texttt{Time:} lines: once 64
tokens remain and no \texttt{Answer:} line has started, the decoder begins one. By itself the reserve
changes S-EMBER mIoU by only $+0.1$ and repairs only three truncated answers.
When the answer is a single interval, PCD has little to correct. With the one-line grounding prompt it changes mIoU
by $-0.1$ (\cref{sec:transfer}), in runs without the reserve. For the official direct-answer prompt we do not
yet have a clean PCD-only comparison on S-EMBER. On NExT-GQA, PCD under that prompt lowers the scored grounding,
because more times are written as minutes:seconds (\cref{sec:res-contrast}).

\subsection{BOLT: question-guided frames}
\label{sec:retrieval}

When the window is much longer than $F$ frames, the choice of frames to encode can matter as much as the choice of
tokens to keep. Our front end samples frames as in BOLT~\cite{bolt}. We embed each second of video once with
SigLIP~2~\cite{siglip2}, which gives a causal 1\,fps index. The relevance $z_j$ of second $j \leq q$ is its z-scored
cosine similarity with the question, and the $F$ frames are taken at evenly spaced quantiles of
\begin{equation}
  p(j) \;=\; \tfrac{1}{2}\cdot\tfrac{1}{q} \;+\; \tfrac{1}{2}\operatorname{softmax}(z)_j ,
\end{equation}
so half of the frames remain uniform. We skip BOLT when uniform frames would be less than 1\,s apart. Once the
question chooses the tokens, BOLT adds only $+0.8$ to $+0.9$ mIoU on S-EMBER, QVHighlights
and ActivityNet. The runs in the main tables include it.
